\begin{afterword}
\summaryhead

The post starts from the author's reductionist thesis: minds use a map with many levels, such as heat and kinetic energy, while reality seems to have only one, the quarks. Carnot could study heat without its kinetic basis, but reality does not seem to simplify its own calculations. So in what sense is ``I am wearing socks'' true, when there are only quarks?

The author takes truth in Tarski's sense, as a comparison between a belief and reality, tested by experiments, which reach outside one's head, not by introspection. No one can compare a neural pattern to a snowflake quark by quark. The post's answer is that a word such as ``snow'' is a promissory note: a pointer to whatever simple boundary an ideal interpreter would draw around the unseen causes of our experience, as Putnam's Twin Earth suggests. Our beliefs about our beliefs are promissory notes too, which the author calls self-consistent rather than circular. So saying that there are no fundamental hands is not saying that there are no hands: higher-level things exist implicitly in the single level of physics, and the map can be true without naming any quark.

\responsehead

The post states two positions clearly and defends them sensibly. The first is an account of reference: a word like ``snow'' points to whatever physical kind lies behind the things we call snow, before anyone knows what that kind is. The second is the view that reducing a thing is not eliminating it. Both are standard positions in philosophy, and the post presents the first, with ``It seems to me,'' as its own suggestion.

The account of reference is close to Putnam's own conclusion. The post credits Putnam's Twin Earth and ``the subsequent philosophical debate''; a reader should know that its own view is close to the conclusion Putnam drew. In Putnam's story, the liquid that Twin Earthers call ``water'' is not water, because the word's reference is fixed by the stuff around the speakers, whatever its structure. That view, semantic externalism, is Putnam's and Kripke's, and the post's promissory note is a version of it.

The distinction the post ends on, between having no fundamental hands and having no hands, is the one the Stanford Encyclopedia describes as the norm: ``most reductive views are realist,'' and elimination is reserved for cases like phlogiston. The post's puzzle about socks arises from its own thesis, and it resolves it the way most reductionists already do.

The discussion of truth assumes more than it states. Tarski's sentence is called ``the classic definition'' of truth and then used to introduce truth as a comparison between belief and reality. As the notes on earlier posts explain, the sentence is accepted by rival theories of truth too, so it does not decide this. The ``postmodernist folks'' are answered without being named, and the harder form of their question, that reading an experiment depends on prior beliefs, is set aside for another post.

The circle at the end, promissory notes about promissory notes, is named ``self-consistent'' and ``reflective coherence'' rather than defended here. Philosophers have defended such circles for basic forms of reasoning, and the author developed the defence in ``Where Recursive Justification Hits Bottom''.

In short: a clear statement of two standard positions, semantic externalism about words like ``snow'' and reduction without elimination, with the defence of its circularity left to a later post.
\end{afterword}
